\scanpage{87}
% archive sheet number: 86
% block text | lines 1-5
Elle est linéaire \unsure{p.r.|par r.} à $O^\Gamma$. Or $O^\Gamma$ est un anneau $\frac12$ local noethérien, et $O$ est \inserted{module} de type fini sur $O^\Gamma$ (bien connu...), \unsure{enfin|alors} la topologie naturelle sur $A$ (donc $A^\Gamma$) est la $=$, qu'on le considère comme $O$-\inserted{module} ou comme $O^\Gamma$-module.

% block text | lines 6-10
Le noyau $N$ de l'application $\varphi$ a donc comme adhérence, dans $\widehat{A}$, le noyau $\widehat{N}$ de l'application prolongée par continuité. Or nous avons vu que ce dernier contient un élément \uline{inversible} $u$. \deleted{\unsure{Il en résulte aussitôt que}} \inserted{Soit $v\in N$ tel que}
% block margin note | left of lines 6-10
% slanted
\margin{\unsure{dégager} un lemme : si quel que soit $n$, on a $f\sim g\bmod m^n A$ \deleted{\unsure{démontrer|dégager}} alors $f\sim g$}
% block display | line 11
\[
v\equiv u\quad(\unsure{\mathfrak r|\mathfrak m|\mathfrak b}A)\qquad\text{où }\mathfrak r\text{ est le radical de }O.
\]
% block text | line 12
Alors $\unsure{u|v}$ est aussi inversible, et on a bien
% block display | line 13
\[
g\unsure{\gamma|\sigma}=u^{-1}f(\gamma)\gamma(u).
\]
% block margin note | left of lines 14-15
\margin{$\Vert$}
% block text | lines 14-15
Cela prouve que $H^1(\Gamma,A^*)\longrightarrow H^1(\Gamma_d,A_m^*)$ et $H^1(\Gamma,A^*)-H^1(\Gamma,\widehat{A}^*)$ sont injectifs.

% block text | lines 16-17
Supposons alors que $\underline{O}$ soit un \uline{anneau local} (donc $\Gamma=\Gamma_d$), \deleted{prouvons que} \inserted{\deleted{Considérons le diagramme}}
% block diagram | lines 18-20 cancelled block
\diagram{\begin{tikzcd}[ampersand replacement=\&]
\deleted{H^1(\Gamma,A^*)}\arrow[r]\arrow[d,"\deleted{\varphi}"]\&\deleted{H^1(\Gamma_i,A^*)}\arrow[d,"\deleted{\varphi_i}"]\\
\deleted{H^1(\Gamma,A/m^*)}\arrow[r,dash,"\deleted{\psi}"]\&\deleted{H^1(\Gamma_i,A/m^*)}
\end{tikzcd}}
% block text | left of diagram
\deleted{est injectif}
% block text | lines 21-22 cancelled block
\deleted{Supposons que nous prouvions que $\varphi$ et $\psi$ sont bijectifs, alors il en résulte (\unsure{l'appliq. à} $\Gamma_i$) que $\varphi_i$ l'est}
